\documentclass[12pt,a4paper]{article}
\usepackage{amsmath,amssymb}
\usepackage{booktabs}
\usepackage{geometry}
\usepackage{graphicx}
\usepackage{float}
\usepackage{listings}
\usepackage{xcolor}
\usepackage{algorithm}
\usepackage{algpseudocode}
\usepackage{titlesec}
\usepackage{setspace}
\usepackage{ragged2e}

%----------------------------------------
% Page Setup for A4
%----------------------------------------
\geometry{
  a4paper,
  top=1in,
  bottom=1in,
  left=1in,
  right=1in
}

\setstretch{1.2}

%----------------------------------------
% Section title formatting
%----------------------------------------
\titleformat{\section}
  {\normalfont\bfseries\fontsize{14}{14}\selectfont}
  {\thesection}{1em}{}

%----------------------------------------
% Code listing style
%----------------------------------------
\lstdefinestyle{matlabstyle}{
  language=Matlab,
  basicstyle=\ttfamily\footnotesize,
  keywordstyle=\color{blue}\bfseries,
  commentstyle=\color{green!60!black},
  stringstyle=\color{orange!80!black},
  frame=none,
  numbers=none,
  breaklines=true,
  showstringspaces=false,
  tabsize=2
}

%----------------------------------------
\begin{document}

%----------------------------------------
% Cover Page
%----------------------------------------
\begin{titlepage}
\centering

\vspace*{0.5cm}
\includegraphics[width=0.22\textwidth]{diu_logo.png}\par
\vspace{0.5cm}

{\Large \textbf{Department of Computer Science and Engineering} \par}
\vspace{0.2cm}
{\Large \textbf{Dhaka International University} \par}

\vspace{0.8cm}

{\normalsize
Semester: 12th \par
Lab Report No: 03 \par
Course Code: CSE-420 \par
Course Title: Image Processing Lab \par
Report Title: Pixel Neighbourhood and Distance Measures in MATLAB \par
}

\vspace{1cm}

{\Large \textbf{Submitted By} \par}
\vspace{0.3cm}

{\normalsize
Rasheduzzaman Rakib \par
Reg. No.: CS-D-77-22-120080 \par
Roll: 37 \par
Batch: D-77 \par
Department of Computer Science and Engineering \par
Dhaka International University \par
}

\vspace{1cm}

{\Large \textbf{Submitted To} \par}
\vspace{0.3cm}

{\normalsize
Md. Muksit Ul Islam \par
Assistant Professor \par
Department of Computer Science and Engineering \par
Dhaka International University \par
}

\vfill

{\normalsize Submission Date: 25th September 2026 \par}
{\normalsize Summer 2026 \par}

\end{titlepage}


%----------------------------------------
% Lab Report
%----------------------------------------

\section{Name}

\textbf{Pixel Neighbourhood and Distance Measures in MATLAB}\\[0.3em]
Topics covered: 4-Neighbour, Diagonal Neighbour, 8-Neighbour,
Euclidean Distance, City Block Distance, Chessboard Distance.

%----------------------------------------

\section{Description}

This lab explores pixel connectivity and distance metrics in digital images
using MATLAB. A $5 \times 5$ integer matrix is used as a sample image.
The centre pixel at position $(3,3)$ is selected and its various neighbours
are identified. Three standard distance measures between two pixels are then
computed.

\begin{enumerate}
  \item \textbf{4-Neighbour:} The four pixels directly adjacent to a pixel
        $(x,y)$ along the horizontal and vertical axes: $(x-1,y)$,
        $(x+1,y)$, $(x,y-1)$, and $(x,y+1)$. For pixel $(3,3)$ these are
        the values 8, 18, 12, and 14.

  \item \textbf{Diagonal Neighbour:} The four pixels located at the four
        diagonal corners of $(x,y)$: $(x-1,y-1)$, $(x-1,y+1)$,
        $(x+1,y-1)$, and $(x+1,y+1)$. For pixel $(3,3)$ these are
        7, 9, 17, and 19.

  \item \textbf{8-Neighbour:} The full $3\times3$ block surrounding $(x,y)$,
        combining both 4-neighbours and diagonal neighbours — a total of
        eight surrounding pixels.

  \item \textbf{Euclidean Distance:} The straight-line distance between two
        pixels $(x_1,y_1)$ and $(x_2,y_2)$, computed as
        $D = \sqrt{(x_2-x_1)^2 + (y_2-y_1)^2}$.
        Between $(3,3)$ and $(5,5)$ this gives $D = 2.83$.

  \item \textbf{City Block Distance:} Also called Manhattan distance, it sums
        the absolute differences along each axis:
        $D = |x_2-x_1| + |y_2-y_1|$.
        Between $(3,3)$ and $(5,5)$ this gives $D = 4$.

  \item \textbf{Chessboard Distance:} The maximum of the absolute differences
        along each axis: $D = \max(|x_2-x_1|,\, |y_2-y_1|)$.
        Between $(3,3)$ and $(5,5)$ this gives $D = 2$.
\end{enumerate}

%----------------------------------------

\section{Code}

\begin{lstlisting}[style=matlabstyle]
clc;
clear;
A=[1 2 3 4 5;
6 7 8 9 10;
11 12 13 14 15;
16 17 18 19 20;
21 22 23 24 25];
disp('image matrix');
disp(A);
x=3;
y=3;
p=A(x,y);
fprintf('selected pixel = %d\n',p);

% 1. 4-neighbour
N4={A(x-1,y),A(x+1,y),A(x,y-1),A(x,y+1)};
disp('4-neighbours');
disp(N4);

% 2. diagonal neighbours
ND=[A(x-1,y-1),A(x-1,y+1),A(x+1,y-1),A(x+1,y+1)];
disp('diagonal neighbours');
disp(ND);

% 3. 8-neighbour
N8=A(x-1:x+1,y-1:y+1);
disp('8-neighbours');
disp(N8);

% 4. euclidean distance
x2=5;
y2=5;
D=sqrt((x2-x)^2+(y2-y)^2);
fprintf('euclidean distance = %.2f\n',D);

% 5. city block distance
D=abs(x2-x)+abs(y2-y);
fprintf('city block distance = %d\n',D);

% 6. chessboard distance
D=max(abs(x2-x),abs(y2-y));
fprintf('chessboard distance = %d\n',D);
\end{lstlisting}

%----------------------------------------

\section{Output}

Figure~\ref{fig:output} shows the MATLAB Command Window output produced by
running lab\_03.m.

\begin{figure}[H]
  \centering
  \includegraphics[width=0.82\textwidth]{Figure_1.png}
  \caption{MATLAB Command Window output showing neighbourhood values and
           distance results for the selected pixel.}
  \label{fig:output}
\end{figure}

Figure~\ref{fig:output} confirms the expected results: the 4-neighbours of
pixel $(3,3)$ (value 13) are 8, 18, 12, and 14; the diagonal neighbours are
7, 9, 17, and 19; the full 8-neighbourhood is the $3\times3$ sub-matrix
centred at $(3,3)$. The Euclidean, City Block, and Chessboard distances
between $(3,3)$ and $(5,5)$ are 2.83, 4, and 2 respectively.

\end{document}
